\scanpage{13}
\begin{center}\small Sur les groupes de Galois attachés aux groupes $p$-divisibles\end{center}
\mspage{127}
Le cas $h=3$ n’est possible que si $S$ est de type $\mathbf{E}_6$; mais la dimension de $\mathbf{E}_6$ est bien trop grande pour que ce groupe puisse être plongé dans $\mathbf{SL}(3)$. De même, le cas $h=4$ n’est possible que pour un groupe de type $\mathbf{D}_n$, avec $n\geqslant 4$ (pour $n\leqslant 3$, $\mathbf{D}_n$ est, soit isomorphe à $\mathbf{A}_n$, soit non simple); comme la dimension de $\mathbf{D}_4$ est strictement supérieure à celle de $\mathbf{SL}(4)$, on conclut comme précédemment.

Ceci achève la démonstration du lemme 2, et, avec elle, celles de la prop.~5 et du th.~3.

\textbf{Remarque.} Même lorsqu’on ne fait plus l’hypothèse ($\mathrm{H}^{*}$.3), l’existence dans $\mathrm{S}$ d’un tore $\varPsi$ de dimension $1$, ayant pour caractère un binôme $n_0\sigma^a+n_1\sigma^b$, peut être utilisée pour déterminer (au moins en partie) la structure de $S$.

\begin{center}\textbf{\S~5. Application aux groupes $\boldsymbol{p}$-divisibles}\end{center}

Conservons les notations des paragraphes précédents. Soit $F$ un \textit{groupe $p$-divisible de hauteur $h$} sur l’anneau $A$ des entiers de $K$ (pour tout ce qui concerne les groupes $p$-divisibles, voir [\textit{9}]). Soit $T$ le module de Tate de $F$, et soit $V=T\otimes\mathbf{Q}_p$. On sait ([\textit{9}], n°~2.4) que $\dim. V=h$, et que $\operatorname{Gal}_K$ opère continûment sur $T$ et sur $V$.

Nous aurons besoin de deux des principaux résultats de la théorie de Tate. Tout d’abord la décomposition de $V_C$ ([\textit{9}], \S~4, cor.~2 au th.~3):

\textbf{Proposition 6.} -- \textit{$V_C=V\otimes C$ possède une décomposition de Hodge-Tate de la forme:}
\[
V_C=V_C(0)\oplus V_C(1).
\]
\textit{De plus, la dimension $n_1$ (resp. $n_0$) de $V_C(1)$ (resp. de $V_C(0)$) est égale à la dimension de la composante neutre de $F$ (resp. de son dual $F'$).}

(Rappelons, cf. [\textit{9}], n°~2.2, que la composante neutre $F^0$ de $F$ peut être identifiée à un \textit{groupe formel} sur $A$, au sens de \textsc{Dieudonné}, \textsc{Lazard}, \textsc{Lubin}; par la «dimension» de $F^0$, on entend le «nombre de paramètres» du groupe formel en question.)

Si $K'$ est une extension finie de $K$, d’anneau de valuation $A'$, notons $\operatorname{End}_{A'}(F)$ l’anneau des endomorphismes du groupe $p$-divisible $F\times_A A'$ déduit de $F$ par extension des scalaires de $A$ à $A'$. Soit $\operatorname{End}(F)$ la limite inductive (qui est en fait une réunion) des $\operatorname{End}_{A'}(F)$ lorsque $K'$ parcourt toutes les sous-extensions finies de $\bar K$ (cf. \textsc{Lubin} [\textit{2}] pour le cas de dimension 1).
% archive sheet number: 12
